\section*{Chapter \ref{ch:qm:stochastic-control} --- Stochastic Control}

\begin{solution}{exo:qm:stochastic-control:1}
$0.04/(4 \times 0.04) = 25\%$.
\end{solution}

\begin{solution}{exo:qm:stochastic-control:2}
$0.02 + 0.05 - \tfrac12 \times 3 \times 0.0324 = 2.14\%$ a year, against 3.29\% at the optimum.
\end{solution}

\begin{solution}{exo:qm:stochastic-control:3}
For $\ln w$: $-w(-1/w^2)/(1/w) = 1$. For $w^{1-\gamma}/(1 - \gamma)$: $-w(-\gamma w^{-\gamma-1})/w^{-\gamma} = \gamma$.
\end{solution}

\begin{solution}{exo:qm:stochastic-control:4}
The certainty-equivalent return $c(\pi) = r + \pi(\mu - r) - \tfrac12\gamma\pi^2\sigma^2$ is a parabola with maximum at $\pi^*$ and
second derivative $-\gamma\sigma^2$, so $c(\pi^*) - c(\pi) = \tfrac12\gamma\sigma^2(\pi - \pi^*)^2 = \tfrac12 \times 3 \times 0.0324 \times 0.0856^2 = 3.6$ bp.
\end{solution}

\begin{solution}{exo:qm:stochastic-control:5}
$0.05/0.0324 = 154\%$, borrowing 54\% of wealth. The fund's risk aversion is higher than log utility, the expected
return is uncertain (overbetting costs more than underbetting, One Quant Book~2, chapter~29), and leverage brings
margin and drawdown constraints the model ignores.
\end{solution}

\begin{solution}{exo:qm:stochastic-control:6}
From the proof of \cref{thm:qm:stochastic-control:merton}: $\partial_tV = -(1 - \gamma)\rho V$ and the supremum is $(1 - \gamma)\rho V$,
so the equation holds. For the fund, $\rho = 0.02 + 0.05^2/(2 \times 3 \times 0.0324) = 3.29\%$.
\end{solution}

\begin{solution}{exo:qm:stochastic-control:7}
Yes: 0.315 at every wealth level, against $(\mu - r)/(\gamma\sigma^2) = 0.309$ and a one-period optimum of 0.309; the small gap
comes from interpolating the \omterm{def:qm:stochastic-control:problem}{value function} on the grid.
\end{solution}

\begin{solution}{exo:qm:stochastic-control:8}
Ten years of returns with 18\% volatility estimate the premium with a standard error of $18\%/\sqrt{10} = 5.7\%$, which
moves the \omterm{def:qm:stochastic-control:fraction}{Merton fraction} by $\pm 59$ points: the data cannot distinguish 0\% from 100\%. Shrink the estimate toward a
prior (chapter~14) and size for the uncertainty, not for the point estimate.
\end{solution}

\begin{solution}{pb:qm:stochastic-control:1}
\textbf{1.} $0.05/(3 \times 0.0324) = 51.4\%$.
\textbf{2.} $\rho = 3.29\%$.
\textbf{3.} 3.6 bp a year.
\textbf{4.} 0.02 bp: 52\% is almost exactly optimal.
\textbf{5.} Merton: sell 0.6 point, practically nothing; the rulebook: buy 8 points back to 60\%.
\textbf{6.} 77\% and 31\%.
\textbf{7.} 41\% and 62\%.
\textbf{8.} 154\%: the Kelly, growth-optimal fraction, with borrowing.
\textbf{9.} 0.52 at every wealth level and date; the one-period optimum is 0.517 (continuous: 0.514).
\textbf{10.} \omterm{def:qm:stochastic-control:utility}{CRRA utility} is homothetic: scaling wealth scales utility, so the best fraction is the same at every
wealth.
\textbf{11.} 3.28\%, 3.31\% and 3.24\% (theory for Merton: 3.29\%).
\textbf{12.} 50\%, 68\% and 81\%.
\textbf{13.} 3.3 bp a year.
\textbf{14.} 5.7\%; $\pi^*$ moves by $0.057/(3 \times 0.0324) = 59$ points either way, from about $-7\%$ to $110\%$.
\textbf{15.} The uncertainty in $\mu$, by far: the committee is debating a few basis points of \omterm{def:qm:stochastic-control:ce}{certainty equivalent}.
\textbf{16.} A band around the target inside which the fund does not trade (chapter~10).
\textbf{17.} The fund would hold a liability-hedging portfolio plus a Merton-like position whose size depends on its
funding surplus; the fraction would depend on wealth.
\textbf{18.} Rebalance mechanically to a target; spend the debate on the inputs (expected return, risk aversion) and
on costs, not on timing.
\textbf{19.} Named result: \emph{the rebalancing committee}: the \omterm{def:qm:stochastic-control:fraction}{Merton fraction} is 51.4\%; holding the rulebook's
60\% costs 3.6 bp a year of certainty-equivalent return, never rebalancing about 3.3 bp, and one standard error in
the expected return moves the fraction by 59 points.
\textbf{20.} Rebalance back to the constant fraction, since nothing about the optimum changed with the price.
\end{solution}

\begin{solution}{iq:qm:stochastic-control:1}
$\pi^* = (\mu - r)/(\gamma\sigma^2)$: the expected excess return (more reward, more stock), risk aversion (more aversion,
less), and variance (more risk, less): a mean-variance trade-off per unit of wealth.
\iqlookfor{the formula and its comparative statics.}
\end{solution}

\begin{solution}{iq:qm:stochastic-control:2}
$V(t, x) = \sup_u\E[\int_t^{t+h}f\,ds + V(t + h, X_{t+h})]$; expand $V(t + h, X_{t+h})$ by Itô, take expectations, subtract $V(t,
x)$, divide by $h$ and let $h \to 0$: $\partial_tV + \sup_u\{\mathcal L^uV + f\} = 0$.
\iqlookfor{Itô inside the \omterm{def:qm:stochastic-control:dp}{dynamic programming} principle.}
\end{solution}

\begin{solution}{iq:qm:stochastic-control:3}
The optimal fraction is constant; a fall lowers the equity share below it, so restoring it means buying. Nothing
in the model's opportunity set changed with the price: returns are independent and identically distributed.
\iqlookfor{constant fraction plus i.i.d. returns.}
\end{solution}

\begin{solution}{iq:qm:stochastic-control:4}
Kelly maximises expected log wealth: Merton with $\gamma = 1$, giving $(\mu - r)/\sigma^2$. More risk-averse investors hold a
fraction $1/\gamma$ of it: fractional Kelly is Merton with $\gamma > 1$.
\iqlookfor{log utility and fractional Kelly.}
\end{solution}

\begin{solution}{iq:qm:stochastic-control:5}
Whether the objective is flat there (then the choice barely matters), the control grid is too coarse, the
quadrature too crude, or the interpolation of the \omterm{def:qm:stochastic-control:problem}{value function} too rough; refine each and see whether the policy
converges.
\iqlookfor{flatness versus numerical error, and a convergence check.}
\end{solution}

\begin{solution}{iq:qm:stochastic-control:6}
When the optimal control switches between extremes, or at a stopping boundary, or with degenerate diffusion: the
\omterm{def:qm:stochastic-control:problem}{value function} has kinks. It is then the unique \omterm{def:qm:stochastic-control:viscosity}{viscosity solution}: at every point, smooth functions touching it
from above or below satisfy the equation's inequalities, which is also the notion under which grid schemes converge.
\iqlookfor{test functions touching from above and below.}
\end{solution}
